% =====================================================================
% Section 5: Experiments
% =====================================================================
%
% Status: Draft Section 5 Experiments v1 (2026-05-16)
% Data sources:
%   - Main graph baselines:
%     run_logs/baseline_qwen32b_graph_top200_gpt4omini_full_20260512
%     run_logs/hipporag_qwen32b_valid_graph_top200_full1000_20260513_r2
%   - NeocorRAG:
%     run_logs/neocorrag_aligned_k3_all32_full1000_20260514
%     run_logs/neocorrag_reader_gpt4omini_full1000_20260516
%   - EvidenceFlow:
%     run_logs/all32_etv3_etv4_pcec_gpt4omini_none_full1000_20260514
%   - Ablations:
%     run_logs/etv4_ablation_relation_coverage_all32_gpt4omini_full1000_20260516
%   - Dense-entry diagnostic:
%     run_logs/etv4_full1000_optimized_equivalence_20260511/reports/etv4_ablation6_dense_only_same_entry_full1000.json
%     run_logs/dense_entry_gpt4omini_reader_full1000_20260516
%
% =====================================================================

\section{Experiments}
\label{sec:experiments}

\subsection{Experimental Setup}
\label{sec:experiments:setup}

\paragraph{Datasets.}
We evaluate on HotpotQA, 2WikiMultiHopQA, and MuSiQue, three standard
multi-hop question answering benchmarks. For each benchmark, we use a
1,000-query evaluation split with gold supporting passages. This setting
allows us to report both retrieval quality and downstream answer quality
under the same reader protocol.

\paragraph{Baselines.}
We compare \methodname{} with graph-based and chain-oriented retrieval
baselines. HippoRAG represents OpenIE-based graph retrieval over a
passage--phrase graph. PropRAG represents proposition-level graph
retrieval with path search over proposition units. NeocorRAG represents
a query-time chain-mining approach that augments the reader context with
mined entity-relation evidence chains. We also include the dense entry
ranking from the same NV-Embed-v2 seed retriever to isolate the effect
of graph-based evidence flow beyond the initial semantic ranking. For
methods that require language-model-assisted indexing, chain extraction,
or question-side analysis, we use Qwen3-32B with reasoning traces
disabled in the output. For the final answer generation step, all methods
use the same GPT-4o-mini reader.

\paragraph{Metrics.}
Retrieval is evaluated by supporting-passage recall at the reader cutoff,
denoted R@5. End-to-end QA is evaluated by exact match (EM) and token F1
after the reader consumes the top five retrieved passages. For ablations,
we additionally report All@5, the fraction of questions for which all
gold supporting passages appear in the reader prefix. All scores are
reported as percentages.

\paragraph{Implementation details.}
All retrieval systems produce ranked passage lists and are evaluated with
the same top-five reader context. All methods are evaluated on the same
1,000 questions and corpus snapshots for each benchmark. HippoRAG and
PropRAG retrieve from graph candidate pools capped at 200 passages.
\methodname{} builds a top-100 native candidate pool from the same
top-200 frontier, performs query-local propagation over its
document-grounded OpenIE graph, and passes the resulting ranked list to
the same reader. Dense embeddings are computed with NV-Embed-v2. The
reader uses temperature 0, and EM/F1 are computed with the same answer
normalization across methods. Prompts and hyperparameters for
question-side analysis and reader generation are provided in the
Appendix. Unless otherwise stated, \methodname{} uses relation-grounded
expansion followed by coverage-aware reranking.

\subsection{Main Results}
\label{sec:experiments:main}

\begin{table*}[t]
\centering
\small
\begin{tabular}{lccc ccc ccc}
\toprule
& \multicolumn{3}{c}{HotpotQA} &
\multicolumn{3}{c}{2WikiMultiHopQA} &
\multicolumn{3}{c}{MuSiQue} \\
\cmidrule(lr){2-4}\cmidrule(lr){5-7}\cmidrule(lr){8-10}
Method & R@5 & EM & F1 & R@5 & EM & F1 & R@5 & EM & F1 \\
\midrule
Dense entry & 93.4 & 59.5 & 73.0 & 75.9 & 55.0 & 60.5 & 68.3 & 35.2 & 46.2 \\
HippoRAG & 92.6 & 59.3 & 72.5 & 82.8 & 56.7 & 63.4 & 71.0 & 34.9 & 46.5 \\
PropRAG & 95.1 & 61.8 & 75.1 & 90.9 & 61.1 & 69.1 & \textbf{74.2} & 36.5 & 47.6 \\
NeocorRAG & 95.0 & 62.1 & 74.8 & 88.9 & 59.8 & 66.5 & 73.1 & 37.1 & 48.1 \\
\methodname{} & \textbf{96.5} & \textbf{62.8} & \textbf{75.6}
& \textbf{96.2} & \textbf{65.9} & \textbf{74.7}
& 73.5 & \textbf{37.2} & \textbf{48.9} \\
\bottomrule
\end{tabular}
\caption{Main results under a unified reader setting. R@5 measures
supporting-passage recall in the top-five reader context, and EM/F1 are
computed from the GPT-4o-mini answer generated from that context.}
\label{tab:main-results}
\end{table*}

Table~\ref{tab:main-results} shows that \methodname{} gives the strongest
average retrieval and QA performance among the evaluated methods. The
clearest gain appears on 2WikiMultiHopQA. Compared with the dense entry
ranking, \methodname{} improves R@5 from 75.9 to 96.2 and F1 from 60.5
to 74.7. Compared with the strongest graph baseline on this dataset,
PropRAG, \methodname{} improves R@5 from 90.9 to 96.2 and F1 from 69.1
to 74.7. This is the setting where recovering bridge evidence is
especially important, and the gain suggests that query-local fact
propagation and enhancement are useful beyond independent semantic,
proposition, or passage ranking. On HotpotQA, the baselines are already
strong, but \methodname{} still gives the best R@5, EM, and F1. On
MuSiQue, PropRAG obtains the highest R@5, while \methodname{} obtains
the best EM and F1. This indicates that higher supporting-passage recall
does not always translate linearly into better answer generation,
especially when questions require longer or more fragmented evidence
combinations.

NeocorRAG provides a strong chain-oriented baseline under the same reader.
Its saved reader contexts contain 3.3 to 3.6 passages on average. This
creates a different precision--coverage trade-off, so we report both
retrieval and QA metrics under the same reader. \methodname{} improves
average F1 from 63.1 to 66.4 and average R@5 from 85.7 to 88.7 over
NeocorRAG. The result supports the central claim that document-grounded
graph propagation can be converted into an effective evidence list
without mining new evidence chains from retrieved text at query time.

\subsection{Ablation Study}
\label{sec:experiments:ablation}

\begin{table*}[t]
\centering
\small
\begin{tabular}{lccc ccc ccc}
\toprule
& \multicolumn{3}{c}{HotpotQA} &
\multicolumn{3}{c}{2WikiMultiHopQA} &
\multicolumn{3}{c}{MuSiQue} \\
\cmidrule(lr){2-4}\cmidrule(lr){5-7}\cmidrule(lr){8-10}
Variant & R@5 & F1 & All@5 & R@5 & F1 & All@5 & R@5 & F1 & All@5 \\
\midrule
Full \methodname{} & \textbf{96.5} & \textbf{75.6} & \textbf{93.3}
& \textbf{96.2} & \textbf{74.7} & \textbf{89.5}
& \textbf{73.5} & 48.9 & \textbf{45.3} \\
Dense entry only & 93.4 & 73.0 & 87.1
& 75.9 & 60.5 & 49.3
& 68.3 & 46.2 & 37.2 \\
w/o relation expansion & 95.2 & 74.8 & 90.9
& 92.8 & 71.1 & 81.1
& 73.0 & \textbf{49.0} & 43.9 \\
w/o coverage reranking & 95.0 & 74.8 & 90.5
& 93.5 & 72.8 & 82.6
& 71.8 & 48.3 & 42.2 \\
\bottomrule
\end{tabular}
\caption{Ablation results for the retrieval framework and enhancement components.
All@5 measures whether the top-five reader context contains all gold
supporting passages.}
\label{tab:ablation}
\end{table*}

Table~\ref{tab:ablation} isolates the retrieval framework and the two
enhancement components in \methodname{}. We use one fixed configuration
across datasets rather than selecting dataset-specific ablation variants.
Replacing
\methodname{} with the dense entry ranking causes the largest drop,
especially on 2WikiMultiHopQA, where F1 decreases by 14.2 points and
All@5 decreases by 40.2 points. This shows that the graph-based evidence
flow is the core source of the retrieval gain, while the enhancement
components refine the final reader prefix.

Removing relation-grounded expansion hurts 2WikiMultiHopQA most clearly,
reducing R@5 by 3.4 points, F1 by 3.6 points, and All@5 by 8.4 points.
This is consistent with the role of expansion in admitting bridge and
tail-entity passages that receive insufficient mass from local
propagation alone. The effect is smaller on HotpotQA and mixed on
MuSiQue. On MuSiQue, the 0.1 F1 difference is small, while the full
model still improves R@5 and All@5. We therefore treat the MuSiQue result
as a retrieval-side gain rather than as strong evidence of a QA gain from
relation expansion.

Removing coverage-aware reranking produces a more uniform retrieval
drop. R@5 decreases on all three datasets, and All@5 decreases by 2.8,
6.9, and 3.1 points on HotpotQA, 2WikiMultiHopQA, and MuSiQue,
respectively. This supports the hypothesis that graph propagation alone
can leave the top of the ranked list concentrated around high-mass
neighborhoods. The noisy-OR coverage objective improves the final reader
prefix by preferring candidates that cover question-side requirements not
already represented by the selected passages.

\subsection{Diagnostic Analysis}
\label{sec:experiments:analysis}

\begin{table*}[t]
\centering
\small
\setlength{\tabcolsep}{4pt}
\begin{tabular}{lccc ccc c}
\toprule
Dataset & \multicolumn{3}{c}{R@5} &
\multicolumn{3}{c}{All@5} &
Support gain \\
\cmidrule(lr){2-4}\cmidrule(lr){5-7}
& Dense entry & \methodname{} & $\Delta$ &
Dense entry & \methodname{} & $\Delta$ & \\
\midrule
HotpotQA & 93.4 & 96.5 & +3.1 & 87.1 & 93.3 & +6.2 & 7.2 \\
2WikiMultiHopQA & 75.9 & 96.2 & +20.4 & 49.3 & 89.5 & +40.2 & 45.1 \\
MuSiQue & 68.3 & 73.5 & +5.3 & 37.2 & 45.3 & +8.1 & 19.9 \\
\bottomrule
\end{tabular}
\caption{Retrieval-side diagnostic comparing the dense entry ranking
with \methodname{} under the same top-five reader prefix. Support gain is
the fraction of queries whose number of retrieved gold supporting
passages increases in the top-five list.}
\label{tab:flow-diagnostic}
\end{table*}

\paragraph{Local evidence flow.}
Table~\ref{tab:flow-diagnostic} directly compares the dense entry
ranking with \methodname{} at the retrieval side. The largest gain
appears on 2WikiMultiHopQA, where R@5 improves by 20.4 points and All@5
improves by 40.2 points. This confirms that the core propagation step is
not merely reordering dense neighbors. It changes the top-five evidence
set for many questions and often increases the number of retrieved gold
supporting passages.

\paragraph{Effect of relation-grounded expansion.}
The largest relation-expansion effect in Table~\ref{tab:ablation} also
appears on 2WikiMultiHopQA, where the task often requires connecting two
entities through an intermediate document. Without relation expansion,
the model still performs local fact-level propagation, but passages
farther from the query anchors are less likely to enter the reader
prefix. The drop in All@5 shows that the missing passages are not merely
alternative evidence but often part of the complete support set required
by the dataset.

\paragraph{Requirement coverage.}
The coverage ablation confirms that selecting the top local-propagation
scores is not enough. Even when the candidate pool contains relevant
passages, the final top-five list can overrepresent one entity
neighborhood. Coverage-aware reranking gives the largest All@5 gain on
2WikiMultiHopQA and improves All@5 on all datasets, which directly
matches the failure mode identified in the introduction.

\paragraph{Retrieval-to-QA conversion.}
The results also show that retrieval metrics and answer metrics are
related but not identical. MuSiQue provides the clearest example.
PropRAG has the highest R@5, while \methodname{} has the highest F1.
This suggests that a useful multi-hop retriever should not only recover
individual gold passages, but also place a connected and covering
evidence set into the reader prefix. This observation motivates reporting
both retrieval and QA metrics throughout the experiments.
